The custom solver described in \autoref{sec:boolean-implementation} is capable of solving the ramp loss model.
This means there is another option, namely the model type, which is either ``Boolean'' or ``Threshold''.
The model type ``Threshold'' solves both the threshold model (\autoref{ch:threshold-model})
and the ramp loss model, as the threshold model is a special case of the ramp loss model,
precisely the case $\beta = 0.$

The implementation for both model types share relevant parts of the codebase.
Notably, both model types use the same basis for the heuristic pricer.
The key difference is how the reduced cost is computed inside the search tree.
We will now discuss how this is done for the ramp loss model.
Let $k_\text{parent}, k_\text{child}$ be the rules corresponding to some parent and child node, respectively.
Then the reduced cost of the child node is given by
\begin{equation}
        \hat \rho_{k_\text{child}} =  \sum_{i \in \cN} \gamma_{ik_\text{child}}
    -\sum_{i \in \cP}  \mu_i \gamma_{ik_\text{child}} +\lambda  c_{k_\text{child}}.  \label{eq:RP_redcost_node}
\end{equation}
We now wish to express $\gamma_{ik_\text{child}}$ in terms of $\gamma_{ik_\text{parent}}$.
Ideally, we can then store $\gamma$-values cleverly in the search tree, and compute
$\hat \rho_{k_\text{child}}$ from $\hat \rho_{k_\text{parent}}$ with relatively little effort.
Denote by $l$ the clause that is added in $k_\text{child}$ with respect to $k_\text{parent}.$
Recall the definition of $\gamma_{ik}$ in \eqref{eq:f_min}.
Evidently,
\[
  \gamma_{ik_\text{child}} \leq \gamma_{ik_\text{parent}}.
\]
Precisely, observe that
\begin{equation}\label{eq:RP_gamma_diff}
    \gamma_{ik_\text{child}} = \gamma_{ik_\text{parent}} - (\gamma_{ik_\text{parent}} - \gamma_{il})\Ind\{\gamma_{il} < \gamma_{ik_\text{parent}}\}.
\end{equation}
For convenience, we denote
\begin{equation}\label{eq:RP_delta_gamma}
    \Delta \gamma_{ik} = (\gamma_{ik_\text{parent}} - \gamma_{il})\Ind\{\gamma_{il} < \gamma_{ik_\text{parent}}\}
\end{equation}
It follows that
\begin{equation}
    \begin{split}
        \hat \rho_{k_\text{child}} = \; & \hat{\rho}_{k_\text{parent}} - \sum_{i \in \cN} \Delta \gamma_{ik_\text{child}}
        + \sum_{i \in \P} \Delta \mu_i \gamma_{ik_\text{child}}
        + \lambda.
    \end{split}
    \label{eq:RP_redcost_node_deriv}
\end{equation}
There are two important cases in which $\Delta \gamma_{ik} = 0$.
First, if $\gamma_{ik_\text{parent}} = 0,$ as $\gamma_{il}$ cannot be lower.
Second, if $\gamma_l = 2,$ as $\gamma_{ik_\text{parent}}$ cannot be higher.
Define
\begin{equation}
    \overline \I_k \coloneqq \{i \in \I \fw \gamma_{ik} > 0\} \tquad k \in \K.\label{eq:ramp_I_k}
\end{equation}
In other words, $\overline \I_k$ is the set of points that are covered by rule $k$ or are in the margin.
Define $\overline \N$ and $\overline \P$ analogously.
Then, instead of summing over $\N$ and $\P$, we may sum over
\[
    \overline \N \cap \{i \in \N : \gamma_{il} < 2\} \quad \text{and} \quad \overline \P \cap \{i \in \P : \gamma_{il} < 2\},
\]
respectively.
Observe that
\[
    \overline \I_{k_\text{child}} \subseteq \overline \I_{k_\text{parent}}.
\]
and
\[
    \overline \I_{k_\text{child}} = \overline \I_{k_\text{parent}} \setminus \{i \in \I \fw \gamma_{il} = 0\}.
\]
We store $\overline \I_k$ for each node $k$ in the search tree.
Furthermore, we do not store every $\gamma_{ik}$, as this would require many variables in each node of the search tree.
Instead, we store the key-value pair $(i, \gamma_{ik_\text{child}})$ at node $k_\text{child}$ if $\gamma_{ik_\text{child}}$ changed with respect to the parent.
Hence, finding the value of $\gamma_{ik_\text{parent}}$ requires traversing up the search tree until there is a hit.

Finally, we note that allowing a TF to appear at most once in a rule can speed up heuristic pricing.
When this restriction is enforced, the number of possible children of each node is reduced.
This is relatively easy to implement for the threshold model and the ramp loss model,
but is less convenient to implement for the boolean model, as it would require additional input to specify which binarized features
belong to the same TF\@.